% id: 148-07
% book: 베이직쎈 미적분1 · 09 정적분의 활용
% section: 학교 시험 기출로 실전 감각 UP!
% figure: tikz:fig-148-07
% answer: 
% answer_source: 
% variant_level: 0 (원본 전사)
% 이 파일은 build.mjs 가 items.json 에서 만든다. 고칠 때는 items.json 을 고친다.
\smash{\raisebox{1.5mm}{\dmkichul{베이직쎈 미적분1 148쪽 07번}}}
\noindent\begin{minipage}[t]{0.58\linewidth}\vspace{0pt}\raggedright 오른쪽 그림과 같이 닫힌구간 $[0{,}\mkern3mu\mathopen{}3]$에서 곡선 $y=-x^2+4$와 세 직선 $x=0$, $x=3$, $y=k$로 둘러싸인 두 도형의 넓이가 같을 때, 상수 $k$의 값은? \dmcondfill{$-5<k<4$}{}\end{minipage}\hfill\begin{minipage}[t]{0.4\linewidth}\vspace{0pt}\centering \resizebox{\ifdim\width>\linewidth\linewidth\else\width\fi}{!}{% 148-07(148쪽) — 곡선 $y=-x^2+4$ 와 세 직선 $x=0$, $x=3$, $y=k$ 로 둘러싸인 두 도형에 색칠. 스캔 화질이 낮아 TikZ 로 다시 그림.
% 원문에 있는 것만 옮긴다: 곡선 · 직선 $x=3$·$y=k$ · 색칠한 두 도형 · 라벨 O, 3, 4, k, x, y, $y=k$, $x=3$, $y=-x^2+4$.
% $y=k$ 는 답($k$ 의 값)이 그림에서 읽히지 않도록 원문과 같은 자리(꼭짓점 4 와 $x$축 사이 위쪽)에 두었고 눈금값은 적지 않았다.
\begin{tikzpicture}[x=0.21cm, y=0.25cm, line width=0.6pt]
  \fill[black!12] plot[domain=0:1.51658, samples=41, smooth] (\x,{4-(\x)*(\x)}) -- (0,1.7) -- cycle;
  \fill[black!12] plot[domain=1.51658:3, samples=41, smooth] (\x,{4-(\x)*(\x)}) -- (3,1.7) -- cycle;
  \draw[->, >=stealth, line width=0.7pt] (-4.2,0) -- (7.3,0) node[below=1pt, font=\small] {$x$};
  \draw[->, >=stealth, line width=0.7pt] (0,-5.6) -- (0,5.0) node[left=1pt, font=\small] {$y$};
  \draw[line width=0.5pt] (-2.6,1.7) -- (7.0,1.7);
  \draw[line width=0.5pt] (3,-5.4) -- (3,4.6);
  \draw[domain=-3.05:3.05, samples=121, smooth] plot (\x,{4-(\x)*(\x)});
  \node[anchor=south west, font=\small] at (3.7,1.85) {$y=k$};
  \node[above=0pt, font=\small] at (3,4.6) {$x=3$};
  \node[above left=-3pt and -2pt, font=\small] at (0,4) {$4$};
  \node[below left=-2pt and -2pt, font=\small] at (0,1.7) {$k$};
  \node[below left=-2pt and -3pt, font=\small] at (0,0) {$\pt{O}$};
  \node[below right=-1pt and -3pt, font=\small] at (3,0) {$3$};
  \node[anchor=north, font=\small] at (-0.3,-5.6) {$y=-x^2+4$};
\end{tikzpicture}
}\end{minipage}\par
\dmchoicesauto{\rule[-3ex]{0pt}{8ex}}{$\dfrac{1}{3}$}{$\dfrac{2}{3}$}{$1$}{$\dfrac{4}{3}$}{$\dfrac{5}{3}$}
